% Subseção: extensões desejáveis (-2) e opcionais (-OPC), juntas para caber em 15 páginas.
% Substitui as antigas secoes/F-*-2.tex e secoes/F-*-OPC.tex (20/11/2026), cujo texto está no
% histórico do Git. Números: os mesmos daquelas seções, com o comando e o JSON registrados
% nos arquivos de feature (docs/plano-de-desenvolvimento/features/feature-<ID>.md):
% F-ROTA-2/OPC: medir_rota_extensoes.py e medir_alternativos_linhas.py (20/11);
% F-ESPERA-2: quadro-horarios-bhtrans.json; F-ESPERA-OPC: sensibilidade-bhtrans.json;
% (F-ESPERA-2, F-EFI-2 e F-EFI-OPC regeradas no feed de 18/11 no i5-13600KF em 20/11)
% F-BALD-2: medir_comparacao_baldeacoes.py (17/11); F-BALD-OPC: nao-dominados-2026-11-20.json;
% F-MET-2/OPC: metricas-extensoes-2026-11-20.json; F-EFI-2: f-efi-2-bhtrans.json;
% F-EFI-OPC: f-efi-opc-bhtrans.json; F-RES-2/OPC: falhas-extensoes-2026-11-20.json.
% Donos: F-BALD-2, Rafael; F-ESPERA-2/OPC e F-EFI-2/OPC, Miguel; as demais, Henrique.
\subsection{Extensões desejáveis e opcionais}\label{sec:extensoes}

Cada pergunta ganhou uma extensão desejável e uma opcional, sobre as mesmas visões e o mesmo
núcleo (Tabela~\ref{tab:extensoes}; tempos na Tabela~\ref{tab:desempenho}). Todas foram
medidas no feed de 18/11, exceto a F-BALD-2 (serviço de 04/11).

\begin{table}[t]
\centering
\caption{Extensões: algoritmo, complexidade e um resultado medido no feed real.}
\label{tab:extensoes}
\scriptsize
\setlength{\tabcolsep}{3pt}
\begin{tabularx}{\textwidth}{@{}l>{\raggedright\arraybackslash}p{3.6cm}%
  >{\raggedright\arraybackslash}p{3.1cm}>{\raggedright\arraybackslash}X@{}}
\toprule
Feature & Algoritmo e visão & Complexidade & Resultado medido \\
\midrule
F-ROTA-2 & A*, $h$ = distância$/v_{\max}$; clique no mapa ligado à parada mais próxima;
  $G_E$ & $O((V + E)\log V)$; parada mais próxima $O(V)$ & Mesmo custo do Dijkstra em 200
  pares; expande 99,7\% dos vértices (mediana) \\
F-ROTA-OPC & Yen com desvios em A*; $G_E$ & $O(k\,n(V + E)\log V)$ & $k = 5$, 8 pares: os 5
  caminhos têm a mesma sequência de linhas em todos \\
F-ESPERA-2 & Chegada mais cedo com busca binária nas partidas; Monte Carlo de saídas; $G_T$ &
  $O((V + E)\log E + E\log(k + 1))$ por consulta & Espera de 1,87~min pelo quadro contra
  15,34~min por $H/2$ no mesmo par \\
F-ESPERA-OPC & Intervalo de uma linha multiplicado; dois $G_E$ e Dijkstra por par &
  $O(n(V + E)\log V)$ & Dobrar a linha 561787 leva a espera de 9,44 a 18,89~min; 3 de 100
  pares ficam mais lentos \\
F-BALD-2 & Dijkstra com estado (vértice, embarques), até $k$ baldeações; $G_L$ e $G_E$ &
  $O((k + 2)(V + E)\log((k + 2)V))$ & Em 40 de 200 pares, menos baldeações custam 1,5~min
  (mediana; máximo 14,5) \\
F-BALD-OPC & Fronteira de Pareto tempo $\times$ baldeações, $B + 1$ buscas com limite; $G_E$ &
  $O(B^2(V + E)\log(BV))$ & 170 de 200 pares com mais de um trajeto não dominado \\
F-MET-2 & Proximidade de Wasserman e Faust, $k$ origens; $G_P$ & $O(k(V + E)\log V)$ &
  $k = 50$ em 2,2~s; top 5 na Av.\ Pres.\ Antônio Carlos \\
F-MET-OPC & Propagação de rótulos e modularidade; $G_P$ não dirigido & $O(V + E)$ por
  iteração & 604 comunidades, $Q = 0{,}806$, igual ao NetworkX \\
F-EFI-2 & Isócrona (Dijkstra limitado a $T$) e índice de desvio; $G_E$ &
  $O((V' + E')\log E')$ na parte visitada & 68, 603, 1.912 e 4.121 paradas em 15, 30, 45 e
  60~min \\
F-EFI-OPC & Eficiência e tempo médio por faixa, mesmos pares; $G_E$ por faixa &
  $O(F\,k(V + E)\log E)$ & Tempo médio de 60,0~min no pico da manhã a 66,0~min à noite \\
F-RES-2 & Ranking: $N$ paradas mais centrais e linhas, cada uma falhando sozinha; $G_P$ &
  $O((c + 1)s(V + E)\log V)$ & 30 candidatos em 38,9~s; maior perda 1,11\% \\
F-RES-OPC & Falhas em sequência; remoção dirigida contra aleatória; $G_P$ &
  $O((k + 1)s(V + E)\log V)$ & 10 remoções: 1,73\% (dirigida) contra 0,24\% (aleatória) \\
\bottomrule
\end{tabularx}
\end{table}

\textbf{Trajetos.} O A* devolveu o mesmo custo do Dijkstra, mas quase não poda: 181 trechos com
horário repetido no GTFS ficam no piso de 1~s e levam $v_{\max}$ a 1.258~m/s, e $h$ ignora as
esperas; a mediana foi de 81,0~ms contra 64,1~ms do Dijkstra. Os alternativos do Yen diferiam
só no ponto de troca ou no trecho a pé, então a página agrupa por sequência de linhas.

\textbf{Espera.} Em $G_T$, entre as paradas 14788851 e 14792190 às 08h, o
quadro de horários deu 23,67~min de porta a porta, contra 40,24~min por $H/2$; cada modelo
escolhe seu trajeto. $G_T$ não tem tempo mínimo de transferência: o trajeto do exemplo tem uma troca
com 0~s de espera, o que favorece o quadro na comparação. Na sensibilidade, os 3 pares afetados
ficam em média 1,35~min mais lentos; nos outros 97 o melhor trajeto não depende da linha.

\textbf{Baldeações.} O tempo comparado é o do mais rápido com até $k$ baldeações, não o do
trajeto de $G_L$, cujo desempate não olha o tempo. A comparação levou 212,9~ms por par em média
e 1.015,5~ms no pior caso, pouco acima do 1~s do RNF-05. A fronteira de Pareto é exata porque o
tempo mínimo com até $k$ baldeações não cresce com $k$, e coincidiu com a força bruta do
NetworkX em 20 redes aleatórias; a maior vai de 0 a 4 baldeações em 151,4 a 98,5~min.

\textbf{Centralidades e eficiência.} A proximidade coincidiu com a do NetworkX em 15 grafos
aleatórios, mas a amostrada não traz erro padrão. As comunidades dependem da semente (587 a 608
com cinco sementes) e saem pequenas: a maior tem 1,5\% das paradas. O maior índice de desvio em
12 pares foi 1,55 (13,7~km percorridos para 8,9~km em linha reta). Entre as faixas, as
diferenças são da ordem de um erro padrão (cerca de 5~min): com 12 pares, não as distinguem.

\textbf{Falhas.} O ranking simula só candidatos (10 paradas de maior intermediação e 20 linhas
sorteadas) com as mesmas 30 origens. Lideram três paradas vizinhas (Rua Deputado Salim Nacur e
Rua das Canárias); a linha 9550 é a 5.ª (0,40\%), e a parada de maior intermediação causa só
0,14\%: centralidade e criticidade não coincidem. No ataque dirigido, a intermediação é calculada
uma vez, na rede intacta.
